%"Computer Algebra Application to Involutivity and Group Analysis of
%Differential Equations"

%List of talks and abstracts:

\documentclass[11pt, draft]{article}
\topmargin=0mm
\oddsidemargin=0mm
\textwidth170mm
\textheight230mm
\begin{document}
%\pagestyle{empty}
\language=0
%\hyphenation{}

\begin{center}
{\bf
INVOLUTIVE ALGEBRAIC AND DIFFERENTIAL SYSTEMS}\\
\vspace {0.3cm}
{Vladimir Gerdt}\\
{\it \small Laboratory of Computing techniques and Automation, Joint Institute for Nuclear
Research, 141980 Dubna, Russia e-mail: gerdt@jinr.ru}
\end{center}
\vspace {1.0cm}
Among the properties of systems of analytical partial differential
equations (PDEs) which may be investigated without their explicit
integration there are compatibility and formulation of an initial-value
problem providing existence and uniqueness of the analytical solution.
The classical Cauchy-Kowalevskaya theorem establishes a class of
quasilinear PDEs which admit posing such an initial-value
problem.

The main obstacle in investigating other classes of PDE systems of
some given order $q$ is existence of integrability conditions,
that is, such relations for derivatives of order $\leq q$ which are
differential but not pure algebraic consequences of equations in
the system.  An involutive system of PDEs has all the integrability
conditions incorporated in it. This means that prolongations of the
system do not reveal integrability conditions. Extension of a
system by its integrability conditions is called completion.
The quasilinear systems of Cauchy-Kowalevskaya type form a particular
family of involutive systems.

In this talk we present some recent achievements in development of
constructive methods for completion to involution of polynomial and
linear differential systems. We argue that polynomial involutive
systems give a useful alternative of the reduced Gr\"{o}bner bases
in commutative algebra. The general involutive approach is based on the
concept of involutive monomial division which is defined for a
monomial set. Every particular division provides for each monomial in
the set the self-consistent separation of variables into multiplicative
and non-multiplicative. An involutive basis is constructed by combining of
non-multiplicative prolongations with multiplicative reductions.

In differential algebra, in addition to their role as canonical bases of
differential ideals, linear and quasilinear (orthonomic) involutive systems
allow one to pose of an initial value problem providing uniqueness of
its solution and to determine the structure of arbitrariness in general
analytical solution.

To illustrate this technique by important example we consider the Navier-Stokes
equations in $R^2$ for an incompressible fluid and determine the functional
arbitrariness in their general analytic solution.

As another important application of the involutivity we consider the
Lie symmetry analysis of nonlinear differential equations and show that by
completion to involution of the determining linear PDEs for Lie symmetry
generators one can easy find the size of symmetry group without explicit
integration of the determining equations.
\end{document}
